problem:  
Let \((M,\Delta,g)\) be a compact, connected, equiregular **step‑2 sub‑Riemannian manifold**, with sub‑Riemannian Hamiltonian  
\[
H(x,p)=\frac12\langle p|_{\Delta_x},p|_{\Delta_x}\rangle_g,
\]
and canonical symplectic form \(\omega_{\mathrm{can}}\) on \(T^*M\).  

Assume that \(H\) is **noncommutatively completely integrable in the Mishchenko–Fomenko sense**, that is:  

1. There exists a finite-dimensional Lie algebra \(\mathfrak g\) of smooth functions on \(T^*M\) such that each \(f\in \mathfrak g\) satisfies \(\{H,f\}=0\), and their Hamiltonian vector fields are complete.  
2. The Poisson structure on \(\mathfrak g\) induced by \(\{\cdot,\cdot\}\) has constant rank \(2r\) on an open dense subset of \(T^*M\).  
3. With \(k=\dim \mathfrak g\), the standard dimension constraint holds:
   \[
   \dim T^*M = k+r,
   \]
   so that generic common level sets of all \(f\in\mathfrak g\) are compact invariant tori of dimension \(k-r\).  
4. The corresponding Hamiltonian group \(G=\exp(\mathfrak g)\) acts properly and symplectically on \(T^*M\) preserving both \(\omega_{\mathrm{can}}\) and \(H\), and the regular energy level \(\Sigma=H^{-1}(\tfrac12)\) is foliated by these compact invariant tori on which \(H\) induces quasi‑periodic motion.  

**Conjectural problem:**  
Suppose additionally that at every point \(x\in M\), the restriction of the natural map \(T_x^*M\to \mathfrak g^*\) (coming from the momentum functions \(J=(f_1,\dots,f_k)\)) has image of constant rank, and that the corresponding distribution on \(M\) generated by the projected vector fields of \(\mathfrak g\) equals \(\Delta\).  

Does it then follow that \((M,\Delta,g)\) arises from a **Riemannian metric \(g_R\)** on \(M\) for which all normal sub‑Riemannian geodesics of \((\Delta,g)\) are, after reparametrization, Riemannian geodesics of \((M,g_R)\)?  
Equivalently, does full noncommutative (Mishchenko–Fomenko) integrability force the sub‑Riemannian geodesic flow to be the **restriction of a Riemannian geodesic Hamiltonian** to a coisotropic submanifold invariant under the \(G\)-action?  

Why is it a "good" problem:  
- **Addresses a precise open boundary:** The problem isolates the tension between maximal Hamiltonian integrability and geometric nonholonomy. The added assumption that the projected symmetries of integrability generate the distribution \(\Delta\) ties the phase‑space structure back to the base manifold, directly testing whether maximal regularity in cotangent dynamics implies a hidden Riemannian structure.  

- **Bridge between fields:** It fuses **sub‑Riemannian geometry**, **integrable Hamiltonian systems**, and **symplectic reduction theory**, demanding interaction of group actions, coisotropic structures, and control geometry.  

- **Profound insight:** The step‑2 case—where \(\Delta\) and its Lie brackets span \(TM\)—is the simplest genuinely nonholonomic setting and thus the first natural testing ground for whether full integrability destroys nonholonomy.  

- **Conceptual and simple formulation:** “Does complete integrability make a nonholonomic system Riemannian in disguise?” is a succinct question that conceals profound geometric consequences.  

- **Potential impact:** Either answer carries weight. A positive result would yield a rigidity theorem linking integrability to Riemannian geometry; a counterexample would identify the first instance of a **truly integrable but genuinely non‑Riemannian sub‑Riemannian geometry**, revealing a new class of Hamiltonian structures and expanding the scope of integrability theory in geometric mechanics.